\documentclass[10pt,a4paper]{amsart}
\usepackage[margin=19mm]{geometry}
\usepackage{amsmath,amssymb,mathtools}
\usepackage[scr=rsfs,cal=euler]{mathalfa}
\usepackage{xfrac}
\usepackage[unicode,hidelinks,bookmarks=false]{hyperref}
\newcommand{\ie}{i.e.\ }
\newcommand{\cf}{cf.\ }
\newcommand{\resp}{respectively\ }
\newcommand{\Subsection}{\subsection}
\providecommand{\nobreakdash}{\nobreak\hbox{-}}
\setlength{\emergencystretch}{2em}
\multlinegap0pt
\usepackage{algorithm}\usepackage{algpseudocode}
\usepackage{xcolor}
\usepackage{amssymb}
\usepackage{bm}
\usepackage{mathtools}
\usepackage[shortlabels]{enumitem}
\usepackage{tcolorbox}
\newtheorem{assumption}{Assumption}
\newtheorem{lemma}{Lemma}
\newtheorem{theorem}{Theorem}
\newcommand{\E}{\mathbb{E}}
\newcommand{\R}{\mathbb{R}}
\title{\vspace{-10mm}\textbf{Beyond Bounded Variance: Variance-Reduced Normalized Methods for Nonconvex Optimization under Blum--Gladyshev Noise}}
\author{Antesh Upadhyay, Arda Fazla, Abolfazl Hashemi}
\date{Source: arXiv:2605.15314v1 (CC BY 4.0)}
\begin{document}
\maketitle

\noindent\emph{Source and licence.} \vspace{-10mm}\textbf{Beyond Bounded Variance: Variance-Reduced Normalized Methods for Nonconvex Optimization under Blum--Gladyshev Noise}. Version arXiv:2605.15314v1 (CC BY 4.0), distributed by arXiv under the Creative Commons Attribution 4.0 International licence, http://creativecommons.org/licenses/by/4.0/, as recorded in the arXiv licence metadata for this exact version and shown on the abstract page https://arxiv.org/abs/2605.15314v1. Full licence notice: \texttt{licenses/arxiv-2605.15314-CC-BY-4.0.txt}.\par\smallskip

\noindent\textit{Editorial scope.} Counted unit: the article's theorem thm:app\_nsgdm\_l1\_master (source0.tex lines 891--927). The article's complete original proof of it is delivered with it (source0.tex lines 929--1138, one \texttt{proof} environment of 210 lines). No mathematics is written, repaired or re-derived: the statement and the proof are the article's own lines, in the article's own notation, and no retained line is substituted or re-flowed.  Retained uncounted as the source-local context the proof needs: lines 297--300; lines 301--304; lines 340--346; lines 727--732; lines 738--757; lines 786--801; lines 861--867; lines 869--878; lines 881--887.  Nothing was omitted from any retained span, so the package carries no omission record.  No retained line contains a citation, so the wrapper carries no bibliography and no bibliography entry.  No retained line contains a figure environment, an image-inclusion command or drawing code, so no image is delivered and the package declares no asset dependency.  Editorial formatting only, in addition: an AMS \texttt{amsart} preamble that restates the article's own declarations and macros used by the retained lines, namely the environments \texttt{assumption}, \texttt{lemma}, \texttt{theorem} on separate counters, together with the standard packages those lines need.  The retained environments print in this document as Assumption 1, Lemma 1, Theorem 1 in order of appearance, whereas the article prints the same items with the numbering of its own full document.  The complete native source of the article as distributed in the arXiv e-print for this exact version is bundled unchanged as \texttt{sources/200645/source0.tex}.
\begin{assumption}[Lower boundedness of $f$]
    \label{as:lower}
    Function $f$ is bounded from below, i.e., $f^{\inf}>-\infty$.
\end{assumption}

\begin{assumption}[Unbiased $\mathsf{BG}$-0 stochastic oracle]
    \label{as:bg0_oracle}
    For every $x\in\R^d$, the stochastic gradient oracle is unbiased, i.e., $\E_\xi[\nabla f(x;\xi)]=\nabla f(x).$ Moreover, there exist constants $B,G\ge0$ such that $\E_\xi \left\| \nabla f(x;\xi)-\nabla f(x) \right\|^2 \le B^2\|x-x^0\|^2+G^2$.
\end{assumption}

\begin{equation}
\label{eq:nsgdm_updates}
\mathsf{NSGDM:}\qquad
x^{k+1} = x^k-\gamma\frac{v^k}{\|v^k\|};
\qquad
v^{k+1} = (1-\eta)v^k+\eta\nabla f(x^{k+1};\xi^{k+1}),
\end{equation}

\begin{equation}
\label{eq:app_uniform_output_identity}
    \E\|\nabla f(\widehat x)\|
    =
    \frac1T\sum_{k=0}^K g_k.
\end{equation}

\begin{lemma}[Trajectory control by normalization]
\label{lem:app_traj_control}
For the \(\mathsf{NSGDM}\) iterates in~\eqref{eq:nsgdm_updates},
\[
    \|x^{k+1}-x^k\|\le \gamma,
    \qquad
    \|x^k-x^0\|\le k\gamma.
\]
Consequently, under Assumption~\ref{as:bg0_oracle},
\begin{equation}
\label{eq:app_bg0_along_path}
    \E\|\nabla f(x^k;\xi^k)-\nabla f(x^k)\|^2
    \le B^2k^2\gamma^2+G^2.
\end{equation}
In particular, since \(v^0=\nabla f(x^0;\xi^0)\),
\begin{equation}
\label{eq:app_B0_bound}
    b_0\le G.
\end{equation}
\end{lemma}

\begin{lemma}[Direction-error inequality]
\label{lem:app_direction_error}
For any \(a,e\in\R^d\), define
\[
    d=
    \begin{cases}
        \dfrac{a+e}{\|a+e\|}, & a+e\neq0,\\[4pt]
        0, & a+e=0.
    \end{cases}
\]
Then
\begin{equation}
\label{eq:app_direction_error}
    \langle a,d\rangle\ge \|a\|-2\|e\|.
\end{equation}
\end{lemma}

\begin{equation}
\label{eq:app_l1_grad_consequence}
    \|\nabla f(y)-\nabla f(x)\|
    \le
    \bigl(L_0+L_1\|\nabla f(x)\|\bigr)
    e^{L_1\|y-x\|}\|y-x\|,
\end{equation}

\begin{equation}
\label{eq:app_l1_func_consequence}
    f(y)
    \le
    f(x)+\langle\nabla f(x),y-x\rangle
    +\frac12
    \bigl(L_0+L_1\|\nabla f(x)\|\bigr)
    e^{L_1\|y-x\|}
    \|y-x\|^2.
\end{equation}

\begin{equation}
\label{eq:app_grad_growth_l1}
    \|\nabla f(x)\|
    \le
    8L_1(f(x)-f^{\inf})+\frac{L_0}{L_1},
    \qquad L_1>0,
\end{equation}

\begin{theorem}[Single-sample master bound for \(L_0\)-smoothness and \(\alpha=1\)]
\label{thm:app_nsgdm_l1_master}
Suppose Assumptions~\ref{as:lower} and~\ref{as:bg0_oracle} hold. Assume
that either \(f\) is \(L_0\)-smooth, or that
\eqref{eq:app_l1_grad_consequence}--\eqref{eq:app_l1_func_consequence} hold.
In the standard smooth case, set \(L_1=0\). If
\begin{equation}
\label{eq:app_l1_master_conditions}
    \gamma L_1\le \frac12,
    \qquad
    \frac{32L_1^2\gamma^2T}{\eta}\le\frac12,
\end{equation}
with both conditions interpreted as vacuous when \(L_1=0\), then
\begin{align}
\E\|\nabla f(\widehat x)\|
&\le
\frac{2\Delta}{\gamma T}
+\frac{4b_0}{\eta T}
+\frac{16L_0\gamma}{\eta}
+2L_0\gamma +4\sqrt{\eta}
\Big(G+\frac{B\gamma K}{2}\Big).
\notag\\
&
\quad
+\frac{64L_1^2\gamma}{\eta}
\Bigg[
4\Delta
+\frac{4\gamma b_0}{\eta}
+\frac{16L_0\gamma^2T}{\eta}
+4\gamma\sqrt{\eta}\,TG
+2B\gamma^2\sqrt{\eta}\,KT
+2L_0\gamma^2T
\Bigg]
\label{eq:app_l1_master_bound}
\end{align}
When \(L_1=0\), the entire \(L_1^2\)-term vanishes.
\end{theorem}

\begin{proof}
Let
\[
    E_k:=v^k-\nabla f(x^k),
    \qquad
    D_k:=\nabla f(x^k)-\nabla f(x^{k+1}),
\]
and
\[
    \zeta_{k+1}
    :=
    \nabla f(x^{k+1};\xi^{k+1})-\nabla f(x^{k+1}).
\]

\paragraph{Step 1: descent.}
Using \eqref{eq:app_l1_func_consequence} with
\(x=x^k\) and \(y=x^{k+1}=x^k-\gamma d^k\), we get
\[
\begin{aligned}
    f(x^{k+1})
    &\le
    f(x^k)-\gamma\langle\nabla f(x^k),d^k\rangle \\
    &\quad
    +\frac{\gamma^2}{2}e^{L_1\gamma}
    \bigl(L_0+L_1\|\nabla f(x^k)\|\bigr).
\end{aligned}
\]
Since \(\gamma L_1\le 1/2\), we have \(e^{L_1\gamma}\le e^{1/2}<2\).
By Lemma~\ref{lem:app_direction_error},
\[
    \langle\nabla f(x^k),d^k\rangle
    \ge
    \|\nabla f(x^k)\|-2\|E_k\|.
\]
Therefore,
\[
\begin{aligned}
    f(x^{k+1})
    &\le
    f(x^k)
    -\gamma\|\nabla f(x^k)\|
    +2\gamma\|E_k\|
    +L_0\gamma^2
    +L_1\gamma^2\|\nabla f(x^k)\|.
\end{aligned}
\]
Using \(\gamma L_1\le1/2\), we obtain
\begin{equation}
\label{eq:app_l1_descent_recursion}
    \Delta^{k+1}
    \le
    \Delta^k-\frac{\gamma}{2}g_k+2\gamma b_k+L_0\gamma^2.
\end{equation}

\paragraph{Step 2: estimator recursion.}
The momentum update gives
\[
    E_{k+1}
    =
    (1-\eta)E_k+(1-\eta)D_k+\eta\zeta_{k+1}.
\]
Unrolling,
\[
    E_k
    =
    (1-\eta)^kE_0
    +\sum_{t=0}^{k-1}(1-\eta)^{k-t}D_t
    +\eta\sum_{t=0}^{k-1}(1-\eta)^{k-1-t}\zeta_{t+1}.
\]
For the drift term, \eqref{eq:app_l1_grad_consequence},
\(e^{L_1\gamma}\le2\), and \eqref{eq:app_grad_growth_l1} imply, for
\(L_1>0\),
\[
\begin{aligned}
    \E\|D_t\|
    &\le
    2\gamma\bigl(L_0+L_1g_t\bigr) \\
    &\le
    2\gamma\left(L_0+L_1
    \left(8L_1\Delta^t+\frac{L_0}{L_1}\right)\right) \\
    &=
    4L_0\gamma+16L_1^2\gamma\Delta^t.
\end{aligned}
\]
For \(L_1=0\), the second term is absent, and the same displayed bound remains
valid with the \(L_1^2\)-term equal to zero.

For the martingale term, conditional unbiasedness and the orthogonality of
martingale differences give
\[
\begin{aligned}
    &\E\left\|
    \eta\sum_{t=0}^{k-1}(1-\eta)^{k-1-t}\zeta_{t+1}
    \right\| \\
    &\quad\le
    \left(
    \eta^2\sum_{t=0}^{k-1}(1-\eta)^{2(k-1-t)}
    \E\|\zeta_{t+1}\|^2
    \right)^{1/2}.
\end{aligned}
\]
By Lemma~\ref{lem:app_traj_control}, for \(t\le k-1\),
\[
    \E\|\zeta_{t+1}\|^2
    \le
    B^2(t+1)^2\gamma^2+G^2
    \le
    B^2k^2\gamma^2+G^2.
\]
Also,
\[
    \sum_{t=0}^{k-1}(1-\eta)^{2(k-1-t)}
    \le
    \frac1\eta.
\]
Hence
\[
    \E\left\|
    \eta\sum_{t=0}^{k-1}(1-\eta)^{k-1-t}\zeta_{t+1}
    \right\|
    \le
    \sqrt{\eta}(G+Bk\gamma).
\]
Combining the drift and martingale estimates yields
\begin{equation}
\label{eq:app_l1_Bk_bound}
    b_k
    \le
    (1-\eta)^kb_0
    +\frac{4L_0\gamma}{\eta}
    +16L_1^2\gamma
    \sum_{t=0}^{k-1}(1-\eta)^{k-t}\Delta^t
    +\sqrt\eta(G+Bk\gamma).
\end{equation}

\paragraph{Step 3: summing the estimator errors.}
Summing \eqref{eq:app_l1_Bk_bound} over \(k=0,\ldots,K\), and using
\[
    \sum_{k=0}^K(1-\eta)^k\le \frac1\eta,
\]
and
\[
    \sum_{k=0}^K\sum_{t=0}^{k-1}(1-\eta)^{k-t}\Delta^t
    \le
    \frac1\eta S_\Delta(K),
\]
we get
\begin{equation}
\label{eq:app_l1_SB_bound}
    S_b(K)
    \le
    \frac{b_0}{\eta}
    +\frac{4L_0\gamma T}{\eta}
    +\frac{16L_1^2\gamma}{\eta}S_\Delta(K)
    +\sqrt\eta
    \left(TG+\frac{B\gamma KT}{2}\right).
\end{equation}

\paragraph{Step 4: controlling the accumulated suboptimality.}
Dropping the negative gradient term in
\eqref{eq:app_l1_descent_recursion}, iterating, and summing gives
\[
    S_\Delta(K)
    \le
    (K+2)\Delta
    +2\gamma T S_b(K)
    +\frac{L_0\gamma^2T(K+2)}{2}.
\]
Substituting \eqref{eq:app_l1_SB_bound} gives
\[
\begin{aligned}
S_\Delta(K)
&\le
(K+2)\Delta
+\frac{2\gamma T b_0}{\eta}
+\frac{8L_0\gamma^2T^2}{\eta}
+\frac{32L_1^2\gamma^2T}{\eta}S_\Delta(K) \\
&\quad
+2\gamma\sqrt\eta T^2G
+B\gamma^2\sqrt\eta KT^2
+\frac{L_0\gamma^2T(K+2)}{2}.
\end{aligned}
\]
By \eqref{eq:app_l1_master_conditions}, the coefficient of
\(S_\Delta(K)\) on the right is at most \(1/2\). Moving it to the left and
using \(K+2\le2T\), we obtain
\begin{equation}
\label{eq:app_l1_Sdelta_avg_bound}
\frac{S_\Delta(K)}{T}
\le
4\Delta
+\frac{4\gamma b_0}{\eta}
+\frac{16L_0\gamma^2T}{\eta}
+4\gamma\sqrt\eta\,TG
+2B\gamma^2\sqrt\eta\,KT
+2L_0\gamma^2T.
\end{equation}

\paragraph{Step 5: final stationarity bound.}
Summing \eqref{eq:app_l1_descent_recursion} gives
\[
    \frac{\gamma}{2}\sum_{k=0}^K g_k
    \le
    \Delta+2\gamma S_b(K)+L_0\gamma^2T.
\]
Dividing by \(\gamma T/2\), using
\eqref{eq:app_uniform_output_identity}, and substituting
\eqref{eq:app_l1_SB_bound} and \eqref{eq:app_l1_Sdelta_avg_bound}, we obtain
\eqref{eq:app_l1_master_bound}.
\end{proof}

\end{document}
